\begin{lemma}[Balanced triples, part (a); Scheide Lemma 3.12(a)]
If $G$ has a globally maximum triple, then its normal-tree height satisfies
$h(G)\ge 3$ and is odd.
\end{lemma}
\begin{proof}
By \noderef{ScheideNormalisation}, at least one normal globally maximum
triple exists.  The possible heights form a nonempty finite set bounded by
$|V(G)|$, so the maximum in \noderef{ScheideGraphHeight} is attained by
some normal data $N$ of height $p=h(G)$.

The closing $(\alpha,\beta)$-path in
\noderef{NormalGlobalMaximumTriple} has distinct multigraph ends $y_1$ and
$y_0$, because tree vertices are injective by \noderef{TreeSequence}.
By \noderef{TwoColorIncidenceGraph}, every edge-node on this incidence path
is an individual multigraph edge identity colored either $\alpha$ or
$\beta$.  Its first edge-node is adjacent to the initial vertex-node
$y_1$, so that multigraph edge is incident with $y_1$.  The definition
\noderef{MissingColors}, together with \noderef{PresentColors}, says that
$\beta$ missing at $y_1$ rules out a $\beta$-colored edge incident with
$y_1$.  Hence the first path edge has color $\alpha$.  Similarly, the last
edge-node is adjacent to the terminal vertex-node $y_0$, and $\alpha$
missing at $y_0$ rules out color $\alpha$ on that incident edge.  Since
only the two displayed colors are allowed, the last path edge has color
$\beta$.

In particular $\alpha\ne\beta$: the first edge is $\alpha$-colored and
incident with $y_1$, so \noderef{PresentColors} makes $\alpha$ present at
$y_1$, whereas $\beta$ is missing there by \noderef{MissingColors}.  The
global field of $N$ is a globally maximum triple; by
\noderef{GlobalMaximumTriple} its first clause is a $k$-triple, and by
\noderef{KTriple} that triple's coloring is proper in the sense of
\noderef{IsProperEdgeColoring}.  Any two consecutive multigraph edges of
the closing path share their intervening vertex.  They are distinct edge
identities because the incidence path is simple, so they are adjacent by
\noderef{AdjacentEdges}.  Properness therefore gives them different
colors.  Since each color is either $\alpha$ or $\beta$ and
$\alpha\ne\beta$, the edge colors alternate.  They start with $\alpha$ and
end with $\beta$, so the positive number of multigraph edges is even.

That number is $p-1$: the stored incidence path has length $2(p-1)$, with
two incidence steps per multigraph edge identity.  Thus $p-1$ is a positive
even integer, so $p$ is odd and $p\ge3$.
\end{proof}
